\title{Temp2TeX Regression Benchmark: Template Fidelity Across Journal Formats}
\author{Alex Example\thanks{Corresponding author: alex.example@example.org} \and Blair Example}
\affiliation{Department of Template Engineering, Example University; alex.example@example.org}
\date{}
\maketitle

\section*{Abstract}
This benchmark manuscript is intentionally structured to exercise journal template behavior. It contains front matter, author notes, an abstract, keywords, multiple heading levels, equations, tables, figures, footnotes, references, declarations, and an appendix. The same body is compiled through the official LaTeX template and the LaTeX package generated by Temp2TeX from the official Word template.

\tempTwoTexKeywords{template conversion; LaTeX; Word template; regression testing; journal formatting}

\section{Introduction}
The regression body verifies whether a converted template preserves the visual contract that a journal author expects from the official source. The body includes paragraphs with citations, inline mathematical symbols, emphasized terms, and cross references. A short footnote appears here to check footnote placement and numbering.\footnote{This footnote is part of the fixed regression body.}

\subsection{Second-level heading}
This section checks heading hierarchy, vertical spacing, paragraph indentation, and line width. It also verifies that a journal template can render normal body text without relying on example-specific macros.

\subsubsection{Third-level heading}
Third-level headings are included because many official Word templates define them even when their LaTeX templates use a different command style.

\section{Methods}
The method section includes an equation and explanatory text. Equation~\ref{eq:regression} should appear with the journal's equation spacing and numbering policy.

\begin{equation}
I = \frac{\sum_{i=1}^{n} w_i x_i}{\sum_{i=1}^{n} w_i}
\label{eq:regression}
\end{equation}

\section{Tables and Figures}
Table~\ref{tab:regression} checks caption placement, table rules, notes, alignment, and merged-cell behavior.

\begin{table}[htbp]
\centering
\caption{Regression table with a note and a merged cell}
\label{tab:regression}
\begin{threeparttable}
\begin{tabular}{@{}p{0.29\linewidth}@{}p{0.29\linewidth}@{}p{0.29\linewidth}@{}}
\toprule
Format element & Expected stress & Evidence target \\
\midrule
\multicolumn{2}{l}{Merged-cell row} & table structure \\
Heading levels & section to subsubsection & text hierarchy \\
Figures & single and paired panels & caption layout \\
\bottomrule
\end{tabular}
\begin{tablenotes}
\small
\item Note: the table intentionally includes a note below the tabular block.
\end{tablenotes}
\end{threeparttable}
\end{table}

\begin{figure}[htbp]
\centering
\fbox{\rule{0pt}{35mm}\rule{0.82\linewidth}{0pt}}
\caption{Single-panel figure placeholder used for layout regression}
\label{fig:single}
\end{figure}

\begin{figure}[htbp]
\centering
\begin{subfigure}{0.46\linewidth}
\centering
\fbox{\rule{0pt}{25mm}\rule{0.82\linewidth}{0pt}}
\caption{Panel A}
\end{subfigure}\hfill
\begin{subfigure}{0.46\linewidth}
\centering
\fbox{\rule{0pt}{25mm}\rule{0.82\linewidth}{0pt}}
\caption{Panel B}
\end{subfigure}
\caption{Two-panel figure placeholder used for subfigure regression}
\label{fig:paired}
\end{figure}

\section{Results}
The result text cites Table~\ref{tab:regression}, Figure~\ref{fig:single}, Figure~\ref{fig:paired}, and Equation~\ref{eq:regression}. The paragraph is long enough to wrap across multiple lines so that the comparison can observe line spacing and text block width. It also includes a parenthetical citation placeholder in plain text (Author et al., 2026).

\section{Discussion}
Discussion text checks body spacing after floats and headings. The converted template should keep the same page frame, header and footer behavior, caption style, and reference placement as the official LaTeX template when both templates receive this same manuscript body.

\section{Conclusion}
This fixed body is not a scientific article. It is a deterministic template stress test for regression comparison.

\section*{Acknowledgements}
The authors thank the template maintainers. No external funding was used.

\section*{Data Availability}
All data in this manuscript are placeholders created for template regression testing.

\begin{thebibliography}{9}
\bibitem[Author and Author(2026)]{regression-ref}
Author A, Author B. Template regression testing for journal manuscript formats. Journal Formatting Methods. 2026;1:1--10.
\end{thebibliography}

\appendix
\section{Appendix Regression Checks}
The appendix verifies appendix heading behavior and numbering after the reference section.
